% Background section

\subsection{Problem Setting}

We consider the abstract problem of finding $u \in U$ such that
\begin{equation}
    \mathcal{R}(u) = 0,
\end{equation}
where $\mathcal{R}: U \to V^*$ is the PDE operator mapping from a trial space $U$ to the dual of a test space $V$.

For example, Poisson's equation with homogeneous Dirichlet boundary conditions:
\begin{equation}
    -\Delta u = f \quad \text{in } \Omega, \qquad u = 0 \quad \text{on } \partial\Omega,
\end{equation}
can be expressed in weak form as: find $u \in H^1_0(\Omega)$ such that
\begin{equation}
    \inner{\nabla u}{\nabla v}_{L^2} = \inner{f}{v}_{L^2} \quad \forall v \in H^1_0(\Omega).
\end{equation}

\subsection{Physics-Informed Neural Networks}

Physics-Informed Neural Networks (PINNs)~\cite{raissi2019physics} approximate the solution using a neural network $u_\theta$ and minimize a loss function based on the PDE residual:
\begin{equation}
    \mathcal{L}_{\text{PINN}}(\theta) = \frac{1}{N} \sum_{i=1}^{N} |\mathcal{R}^s(u_\theta)(x_i)|^2 + \lambda \cdot \text{BC penalty},
\end{equation}
where $\mathcal{R}^s$ denotes the strong-form residual and $\{x_i\}$ are collocation points.

\subsection{Deep Fourier Residual Method}

The DFR method~\cite{taylor2023deep} uses the weak formulation and minimizes the $H^{-1}$ norm of the residual:
\begin{equation}
    \mathcal{L}_{\text{DFR}}(\theta) = \norm{\mathcal{R}^w(u_\theta)}_{H^{-1}}^2.
\end{equation}

For rectangular domains, the $H^{-1}$ norm can be computed efficiently using the Discrete Sine/Cosine Transform:
\begin{equation}
    \norm{R}_{H^{-1}}^2 = \sum_{k} \frac{|\hat{R}_k|^2}{1 + |k|^2},
\end{equation}
where $\hat{R}_k$ are the Fourier coefficients.

\subsubsection{Key Properties}

The DFR method enjoys the following properties:
\begin{proposition}[Error equivalence~\cite{taylor2023deep}]
    For well-posed problems, there exist constants $\gamma, M > 0$ such that
    \begin{equation}
        \frac{1}{M}\norm{\mathcal{R}(u)}_{V^*} \leq \norm{u - u^*}_{U} \leq \frac{1}{\gamma}\norm{\mathcal{R}(u)}_{V^*}.
    \end{equation}
\end{proposition}

This error equivalence ensures that minimizing the DFR loss is equivalent to minimizing the error, providing reliable training dynamics.

\subsection{Limitations of Current Approaches}

While the DFR method addresses several limitations of PINNs, some challenges remain:
\begin{itemize}
    \item Restriction to rectangular domains
    \item Extension to time-dependent problems
    \item Scalability to high dimensions
    \item [Add other limitations relevant to your work]
\end{itemize}

[Continue with background relevant to your specific extension]
